\section{System and Evaluation Design}
\label{sec:threat}

\begin{figure}[t]
  \centering
  % Figure 1: evaluation chain from untrusted source through a generative
% sender to a receiver, gate and store, with the logged outcome rail aligned
% to the stage at which each outcome is recorded. Line style marks how each
% link is evaluated. Vector source; sized for one column.
\begin{tikzpicture}[
  font=\footnotesize,
  >={Stealth[length=1.5mm]},
  stage/.style={draw, rounded corners=2pt, align=center, text width=24mm,
    minimum height=7mm, inner sep=1.5pt, fill=white},
  ctrl/.style={draw, rounded corners=2pt, align=center, text width=12mm,
    minimum height=6mm, inner sep=1pt, font=\scriptsize},
  rung/.style={circle, draw, inner sep=0pt, minimum size=3.2mm,
    font=\tiny},
  lab/.style={font=\scriptsize, inner sep=1pt, xshift=1pt},
  outtxt/.style={font=\scriptsize, anchor=west, inner sep=1pt},
  tie/.style={densely dotted, black!45},
  model/.style={->},
  replay/.style={->, dashed},
  declared/.style={->, densely dotted, thick},
]
  % Stage column.
  \node[stage, fill=black!8] (src) at (2.65,0)
    {Untrusted content\\[-2pt]{\scriptsize marked or unmarked}};
  \node[stage] (snd) at (2.65,-1.15)
    {Sender agent\\[-2pt]{\scriptsize derived message}};
  \node[stage] (rcv) at (2.65,-2.30)
    {Receiver agent\\[-2pt]{\scriptsize structured actions}};
  \node[stage] (mon) at (2.65,-3.45)
    {Reference monitor\\[-2pt]{\scriptsize gate $\widehat P$ (Table~\ref{tab:contract})}};
  \node[stage] (sto) at (2.65,-4.60)
    {Protected store\\[-2pt]{\scriptsize process-local}};
  \node[stage, fill=blue!5] (dwn) at (2.65,-5.75)
    {Downstream agent\\[-2pt]{\scriptsize reads the store}};

  \draw[model] (src) -- node[lab, right] {relay} (snd);
  \draw[model] (snd) -- node[lab, right] {derived text} (rcv);
  \draw[model] (rcv) -- node[lab, right] {attempt} (mon);
  \draw[replay] (mon) -- node[lab, right] {allow} (sto);
  \draw[->, dashed, black!45] (sto) -- node[lab, right] {possible read} (dwn);

  % Control states enter at the stage they govern.
  \node[ctrl] (a) at (0.62,-2.30) {elevation\\[-2pt]$a$};
  \node[ctrl] (b) at (0.62,-3.45) {write grant\\[-2pt]$b$};
  \draw[model] (a) -- (rcv);
  \draw[declared] (b) -- (mon);

  % Outcome rail: each rung sits at the stage where it is logged.
  \def\railx{4.55}
  \node[font=\scriptsize\bfseries, anchor=west, inner sep=0pt]
    at (\railx-0.1,0) {Logged outcome rail};
  \draw[line width=0.9pt, black!35] (\railx,-0.80) -- (\railx,-6.10);
  \foreach \s/\y in {snd/-1.15, rcv/-2.30, mon/-3.45, sto/-4.60, dwn/-5.75}
    \draw[tie] (\s.east) -- (\railx,\y);

  \node[rung] (rT) at (\railx+0.32,-1.15) {T};
  \node[outtxt] at (rT.east) {marker, action in text};
  \node[rung] (r1) at (\railx+0.32,-2.13) {1};
  \node[outtxt] at (r1.east) {attempt or delegation};
  \node[rung] (r2) at (\railx+0.32,-2.49) {2};
  \node[outtxt] at (r2.east) {unauthorized attempt};
  \node[rung] (r3) at (\railx+0.32,-3.45) {3};
  \node[outtxt] at (r3.east) {gate decision};
  \node[rung] (r4) at (\railx+0.32,-4.42) {4};
  \node[outtxt] at (r4.east) {mutation};
  \node[rung, fill=black, text=white] (r5) at (\railx+0.32,-4.78) {5};
  \node[outtxt] at (r5.east) {\textbf{unauthorized mutation}};
  \node[rung] (r6) at (\railx+0.32,-5.75) {6};
  \node[outtxt] at (r6.east) {exposure};
  \foreach \r in {rT,r1,r2,r3,r4,r5,r6}
    \draw[black!35] (\railx,0 |- \r.center) -- (\r.west);

  % Key.
  \node[font=\scriptsize, anchor=north west, align=left, text width=82mm,
    inner sep=0pt] at (-0.05,-6.27)
    {\tikz[baseline=-0.6ex]\draw[model] (0,0) -- (0.5,0); model calls\quad
     \tikz[baseline=-0.6ex]\draw[replay] (0,0) -- (0.5,0); deterministic gate or replay\quad
     \tikz[baseline=-0.6ex]\draw[declared] (0,0) -- (0.5,0); declared fault weighting (\S\ref{sec:covariance})\\
     Relay studies stop at rungs T--1. Rung 5 is primary only in the
     state-change fixture.};
\end{tikzpicture}

  \caption{Evaluation chain. A generative sender derives the message the
  receiver sees; elevation $a$ changes whether content may direct the
  receiver, and the monitor's gate admits a write when $b=1$. Each outcome
  sits at the stage that logs it. No single experiment traverses the whole
  chain: each fixture covers the rungs listed in Table~\ref{tab:contract}.}
  \Description{Top-to-bottom flow from marked or unmarked untrusted content
  through a sender agent that derives a message, a receiver agent, a
  reference monitor, a protected store, and a downstream agent. Elevation a
  enters at the receiver; write grant b enters the monitor as a dotted
  declared-fault link. Solid links are model calls and dashed links are
  deterministic gate or replay. A rail aligns logged outcomes with stages:
  transformation at the sender, attempt and unauthorized attempt at the
  receiver, gate decision at the monitor, mutation and unauthorized mutation
  at the store, and exposure downstream.}
  \label{fig:system}
\end{figure}

\paragraph{Boundary and controls.}
Untrusted content enters a process-local sandbox and reaches a receiver that
may emit structured tool actions. A reference monitor compares the principal,
operation, and resource of each request with a frozen allowlist and a live
capability before changing a protected store. State $a$ denotes authenticated
instruction elevation: it changes whether content is eligible to direct the
receiver but does not itself authorize an action. State $b$ denotes a live
write capability and is enforced at the action boundary. Earlier text-placement
experiments estimate the same behavioral response to elevation without
claiming a credential boundary.

\paragraph{Outcomes and authorization.}
An \emph{attempt} is a structured tool action, rather than free-text assent.
A \emph{broad write attempt} calls the write tool; an \emph{exact protected
request} specifies the protected resource and attack value. A
\emph{mutation} changes stored state, an \emph{unauthorized mutation} does so
without policy permission, and \emph{exposure} means a downstream agent reads
the attack value. We log policy decisions separately from these outcomes.
Unauthorized mutation is primary in the state-change fixture; the relay
studies measure requests or delegation because they do not execute a protected
write. The principal \texttt{receiver} is never entitled to the protected
canary ($P^\star=0$ throughout); fixtures differ in what the gate
$\widehat P$ checks (Table~\ref{tab:contract}). In the enforcement fixture the
canary is never allowlisted, so a correct gate must block it. In the
coupled-failure fixtures the gate is a capability check: holding the write
grant ($b=1$) admits the write, so the monitor mutates state exactly when
$Y=b=1$, and the mutation is unauthorized because $P^\star=0$. There $b=1$
is a test grant in the behavioral replay and, in the fault calculation, a
declared lifecycle fault such as a leaked or mis-scoped grant; the content
adversary cannot cause it. Authorization is therefore exact and needs no
semantic adjudication. No action touches the host filesystem, network, or
real credentials.

\begin{table}[t]
  \centering
  \footnotesize
  \caption{Execution contract per fixture. $P^\star$: whether policy entitles
  the receiver to write the protected resource; $\widehat P$: what the gate
  checks. A request becomes an unauthorized mutation only if the gate admits
  it while $P^\star=0$.}
  \label{tab:contract}
  \setlength{\tabcolsep}{3pt}
  \begin{tabular}{p{.29\columnwidth}p{.36\columnwidth}p{.27\columnwidth}}
    \toprule
    Fixture & Gate $\widehat P$ on protected write & Counted outcome \\
    \midrule
    Relay, propagation (\S\ref{sec:laundering}, \S\ref{sec:propagation-selector}) & None; no write is executed & Request or delegation \\
    Enforcement (\S\ref{sec:defense}) & Exact allowlist; canary never listed (one arm unenforced) & Mutation \\
    Tool visibility (\S\ref{sec:defense}) & Write denied; escalation non-mutating & Exact write or escalation request \\
    Coupled failure (\S\ref{sec:covariance}) & Capability: admitted iff $b=1$ & $Y$ live; mutation $=Yb$ by replay \\
    AgentDojo (\S\ref{sec:agentdojo}) & Official simulated tools & Official predicates \\
    \bottomrule
  \end{tabular}
\end{table}

\paragraph{Adversary.}
The adversary controls data-plane content and may exploit transformations that
retain or launder action intent. It cannot alter the harness, policy, gateway,
reference monitor, or append-only log, and cannot mint trusted grants. Provider
or model-weight compromise is out of scope. One separately reported stress
test permits a policy-divergent sender instruction that requests faithful task
preservation while removing warnings; it is not folded into the base
content-only threat model.

\paragraph{Models, freezing, and inference.}
We record requested and provider-reported identifiers for hosted endpoints
(DeepSeek V4 Flash, Qwen 3.7 Flash 2026-07-15, and MiniMax-M3) at temperature
zero. Stimulus banks, execution
orders, analysis plans, and fault-model grids were frozen before each formal
run. Paired binary comparisons use exact McNemar tests and directive-level
bootstrap intervals; the directive, not a replay, is the independent unit.
Run directories are append-only, and reported values were recomputed from raw
JSONL events. Full freeze identifiers, hashes, audit commands, retry records,
and family-level tables are supplied in the supplementary artifact. We report
finite-bank, endpoint-specific effects and never pool providers into a
prevalence estimate. An additional run exposed a DeepSeek identifier mismatch;
it is not pooled with the earlier endpoint result (\S\ref{sec:propagation-selector}).
The native AgentDojo screen (\S\ref{sec:agentdojo}) used Qwen 3.6 Flash and a
DeepSeek endpoint that reported \path{deepseek-flash}; neither is pooled with
the endpoints above.

\begin{table*}[t]
  \centering
  \footnotesize
  \caption{Map of claim-bearing evaluations. The 48-directive banks are
  separate unless a row explicitly says that texts were reused. D denotes
  DeepSeek V4 Flash and Q denotes Qwen 3.7 Flash; Q3.6 is Qwen 3.6 Flash and DF
  the endpoint reporting \texttt{deepseek-flash}. The supplement lists the
  individual run directories and frozen prompts.}
  \label{tab:experiment-map}
  \setlength{\tabcolsep}{3pt}
  \begin{tabular}{p{.14\textwidth}p{.16\textwidth}p{.17\textwidth}p{.24\textwidth}p{.19\textwidth}}
    \toprule
    Evaluation & Bank and unit & Sender / receiver & Varied factor & Measured outcome \\
    \midrule
    Five forms (\S\ref{sec:laundering}) & 48 directives, five forms & D / D; same archived texts / Q & Visible marker and sender derivation & Protected-write request \\
    Adaptive relay (\S\ref{sec:adaptive-sender}) & Separate 48-directive bank & D / D & Sender instruction and marker & Broad and exact requests \\
    Two receivers (\S\ref{sec:propagation-selector}) & 192 directives; 121 selected & Archived D / D / D & Frozen retention selector & Two consecutive \texttt{delegate} calls \\
    Enforcement (\S\ref{sec:defense}) & Separate 48-directive bank & Q receiver & Monitor policy and capability & Unauthorized mutation \\
    Tool visibility (\S\ref{sec:defense}) & Separate 48-directive bank & Q receiver & Direct-write tool shown or hidden & Direct write or escalation request \\
    Representation (\S\ref{sec:spotlight-relay}) & Separate 48-directive bank & D / D & Prose, datamarking, Base64 & Direct and relayed write request \\
    Credential gate (\S\ref{sec:covariance}) & Separate 48-directive bank & Q receiver & Elevation grant; replay under write grant & Exact protected request \\
    Relay gateway (\S\ref{sec:e4-gateway}) & E2 bank; reused 48 directives & Frozen D outputs / D; Q repeat & Sender policy and elevation & Broad and exact requests \\
    Protocol paths (\S\ref{sec:external}) & Fixed synthetic cases & Pinned MCP / A2A code & Storage root and message derivation & Credential or lineage state \\
    AgentDojo (\S\ref{sec:agentdojo}) & 55 official Slack pairs & Q3.6 / Q3.6; DF / DF & Visible marker and sender derivation & Official task predicates \\
    \bottomrule
  \end{tabular}
\end{table*}

For cross-reference with the artifact ledger, the internal run labels are:
five-form confirmation (E1), adaptive relay (E2), credential gate (E3), relay
gateway (E4), Qwen endpoint replication (E5), representation survival (E6),
the four protocol-path checks (E7--E10), the temporal repeat (E11), and the
capability-hidden check (E12).
